%% notes-tex/010-additive-baselines/sections/8-limits.tex

\section{What this establishes, and what it does not\secstatus{todo}}
\label{sec:limits}

\subsection{Established\secstatus{todo}}

The additive null had never been fit on 010, and now it has, on the same
records, split and label the transformer used, with the ridge penalty chosen on
validation and reported on test.

Every model in the comparison receives the same information, the identity of the
three perturbed genes. That is a reading of the code rather than an assumption:
the encoder is applied at batch size one to a strain-independent input, and the
nine interaction graphs reach the prediction only through the weights that
training left behind.

The transformer's margin over a model that cannot represent interaction is real
and small on the published split.

The published split was measuring screen identity more than trigenic biology, and
this is now measured rather than suspected. A model that never looks at the third
gene reaches 97 percent of the additive null's score, and making query doubles
disjoint drops the additive null from 0.400 to $0.127 \pm 0.033$ and inverts the
ordering of the baselines.

The graphs are load-bearing during training. They carried 99.99 percent of the
loss, their heads recover their graphs, and the matched pair of runs without the
penalty does not train.

The design-space defect has a confirmed cause, a uniform 28-position gene index
shift, and the artifacts it invalidates are enumerated.

\subsection{Not established\secstatus{todo}}

\notesec{The reproduction is faithful but not bit-exact}
The hand-built forward pass matches the recorded metrics to the fourth decimal,
not exactly. The original runs used mixed bf16 precision and this one runs in
fp32, which is the likely source, but that has not been verified by rerunning in
bf16. Conclusions in Section~\ref{sec:paired} rest on correlations at the second
decimal and are not sensitive to a fourth-decimal shift.

\notesec{One seed of the nonlinear baseline enters the paired analysis}
Table~\ref{tab:paired} uses B5 seed 0. The three seeds agree closely on
aggregate metrics, but their record-level agreement with each other was not
measured, so the B5 column should be read as one draw.

\notesec{The transformer has not been measured on the disjoint split}
Table~\ref{tab:disjoint} has no transformer row. Its checkpoints trained on the
random split, so scoring them on held-out query pairs would measure memorization.
Whether the transformer's $+0.04$ margin survives a disjoint split is not
measured, in either direction.

\notesec{Whether the graphs carry information is separate from whether they are
load-bearing}
The runs without the graph penalty fail to train, and the penalty was almost the
entire loss. Both facts are recorded. Neither separates the graphs supplying
biological structure from a large auxiliary term supplying conditioning that the
optimization needs. The control that would separate them, a comparable penalty
against degree-matched random graphs, has not been run.

\notesec{The graph ablation is short and unreplicated in one arm}
The matched pair ran 30 epochs against the 49 to 64 of the reported checkpoints,
and the arm without the penalty produced two numbers and one diverged run rather
than three.

\notesec{Nothing here bears on the architecture's merits elsewhere}
The reduction in Section~\ref{sec:setfunction} is specific to the 010
configuration, where the masking and propagation arms are off and the
perturbation transform is one block. A configuration that switches those on
would have a strain-dependent encoder path and would need its own reading.

\subsection{Next\secstatus{todo}}

\begin{enumerate}
\item Train the transformer on the query-pair-disjoint split and fill the empty
  row of Table~\ref{tab:disjoint}. The migrated build makes this runnable, and it
  is the measurement the rest of this document is now waiting on.
\item Regenerate the panel-12 and panel-24 selections from a correctly indexed
  inference run. The current selections are invalid, not merely uncertain.
\item Run the graph ablation against degree-matched random graphs, to separate
  the graphs carrying structure from a large penalty carrying conditioning.
\item Pin the gene index space with the build rather than reconstructing it from
  a live genome, which is the fix that would have prevented the defect the guard
  now only detects.
\item Report an ensembled ranking for any future selection rather than one
  checkpoint's tail, given that two training runs share 0.39 to 0.47 of their
  top 100.
\item Give experiment builds the manifest that dataset builds already get, so a
  run can name its data version.
\end{enumerate}
